% BEGIN VOL07-EXPANSION VII41
\section{Supplementary worked examples}
\label{sec:vii41-pedagogy}

\begin{example}[Equilateral triangle stiffness]\label{ex:vii41-ped-01}
For one equilateral triangle, each opposite angle is \(\pi/3\), so a boundary-edge cotangent weight is \(\tfrac12\cot(\pi/3)=1/(2\sqrt3)\).  The local stiffness matrix has negative off-diagonals and row sums zero.
\end{example}

\begin{example}[Constant null mode]\label{ex:vii41-ped-02}
Because the hat functions sum to one, \(K\mathbf1=0\).  This discrete identity is the finite-element analogue of the Laplacian annihilating constant functions.
\end{example}

\begin{example}[Obtuse cotangent weight]\label{ex:vii41-ped-03}
If an angle opposite an edge is obtuse, its cotangent is negative.  The assembled cotangent operator can therefore contain negative individual edge weights even though the global finite-element stiffness matrix remains positive semidefinite.
\end{example}

\section{Graded supplementary exercises}
The following exercises add a deliberate progression from concrete calculation to structural proof, hypothesis testing, application, and synthesis.

\subsection*{Standard computations and constructions}

\begin{exercise}[Cotan sixty]\label{exr:vii41-ped-01}
Compute \(\cot(\pi/3)\).
\end{exercise}

\begin{hint}
Use \(\cos/\sin\).
\end{hint}

\begin{solution}
\(1/\sqrt3\).
\end{solution}

\begin{exercise}[Boundary weight]\label{exr:vii41-ped-02}
What is the cotangent weight of an edge opposite a \(60^\circ\) angle in a single boundary triangle?
\end{exercise}

\begin{hint}
Use one half of the cotangent.
\end{hint}

\begin{solution}
\(1/(2\sqrt3)\).
\end{solution}

\begin{exercise}[Row sum]\label{exr:vii41-ped-03}
What is \(K\mathbf1\)?
\end{exercise}

\begin{hint}
Use partition of unity.
\end{hint}

\begin{solution}
Zero.
\end{solution}

\begin{exercise}[Lumped mass]\label{exr:vii41-ped-04}
How much barycentric lumped mass does a triangle of area \(A\) contribute to each vertex?
\end{exercise}

\begin{hint}
Split area equally.
\end{hint}

\begin{solution}
\(A/3\).
\end{solution}

\begin{exercise}[Pointwise Laplacian]\label{exr:vii41-ped-05}
With positive stiffness \(K\) and mass \(M\), what matrix approximates the nonpositive smooth \(\Delta\)?
\end{exercise}

\begin{hint}
Use this volume's sign convention.
\end{hint}

\begin{solution}
\(-M^{-1}K\).
\end{solution}

\subsection*{Proofs}

\begin{exercise}[PSD stiffness]\label{exr:vii41-ped-06}
Why is \(K\) positive semidefinite?
\end{exercise}

\begin{hint}
Use the Dirichlet energy identity.
\end{hint}

\begin{solution}
\(f^TKf=\int\|\nabla f_h\|^2dA\ge0\).
\end{solution}

\begin{exercise}[Symmetry]\label{exr:vii41-ped-07}
Why is the cotangent stiffness matrix symmetric?
\end{exercise}

\begin{hint}
The bilinear Dirichlet form is symmetric.
\end{hint}

\begin{solution}
\(K_{ij}=\int\langle\nabla\varphi_i,\nabla\varphi_j\rangle dA=K_{ji}\).
\end{solution}

\begin{exercise}[Constant harmonic]\label{exr:vii41-ped-08}
Prove every constant vertex function is discrete harmonic in the interior.
\end{exercise}

\begin{hint}
Use row sums zero.
\end{hint}

\begin{solution}
Multiplying \(K\) by a constant vector gives zero.
\end{solution}

\begin{exercise}[Generalized eigenproblem]\label{exr:vii41-ped-09}
Why is \(Kx=\lambda Mx\) preferable to eigenanalysis of \(K\) alone for geometric spectra?
\end{exercise}

\begin{hint}
The mass matrix represents the discrete \(L^2\) inner product.
\end{hint}

\begin{solution}
The generalized problem incorporates both energy and area weighting, approximating the Laplace--Beltrami spectrum rather than the stiffness form alone.
\end{solution}

\subsection*{Counterexamples and hypothesis tests}

\begin{exercise}[Cot weights always positive]\label{exr:vii41-ped-10}
Test the claim on obtuse triangles.
\end{exercise}

\begin{hint}
Cotangent of an obtuse angle is negative.
\end{hint}

\begin{solution}
False.  Individual cotangent weights can be negative on non-Delaunay or obtuse configurations.
\end{solution}

\begin{exercise}[Stiffness equals graph Laplacian]\label{exr:vii41-ped-11}
Test the claim.
\end{exercise}

\begin{hint}
Graph weights need not be cotangent geometry weights, and mass is separate.
\end{hint}

\begin{solution}
False.  The FEM stiffness resembles a weighted graph Laplacian algebraically but its weights come from triangle geometry and the pointwise operator also involves mass.
\end{solution}

\begin{exercise}[Mass irrelevant]\label{exr:vii41-ped-12}
Test the claim for spectral or pointwise Laplacian computations.
\end{exercise}

\begin{hint}
Stiffness is a bilinear form, not the whole operator.
\end{hint}

\begin{solution}
False.  The mass matrix is needed to represent the \(L^2\) metric and convert weak forms to pointwise/generalized operators.
\end{solution}

\subsection*{Applications and investigations}

\begin{exercise}[Smoothing]\label{exr:vii41-ped-13}
Why does the Laplacian support mesh smoothing and diffusion?
\end{exercise}

\begin{hint}
It penalizes Dirichlet energy/high-frequency variation.
\end{hint}

\begin{solution}
Evolution driven by the Laplacian damps oscillatory modes and propagates information according to intrinsic connectivity and geometry.
\end{solution}

\begin{exercise}[Poisson editing]\label{exr:vii41-ped-14}
What linear system underlies many scalar Poisson editing problems on meshes?
\end{exercise}

\begin{hint}
Use stiffness with boundary/gauge constraints.
\end{hint}

\begin{solution}
A constrained system of the form \(Kf=b\), or its mass-weighted counterpart depending on the weak formulation.
\end{solution}

\subsection*{Challenge problems}

\begin{exercise}[Obtuse sign]\label{exr:vii41-ped-15}
If \(\alpha=120^\circ\), what is the sign of \(\cot\alpha\)?
\end{exercise}

\begin{hint}
Cosine is negative, sine positive.
\end{hint}

\begin{solution}
Negative.
\end{solution}

\begin{exercise}[Energy challenge]\label{exr:vii41-ped-16}
If \(f\) is constant on each connected component, what is its discrete Dirichlet energy?
\end{exercise}

\begin{hint}
All edge differences vanish.
\end{hint}

\begin{solution}
Zero; conversely on a connected nondegenerate mesh, zero FEM Dirichlet energy forces the piecewise-linear function to be constant.
\end{solution}

% END VOL07-EXPANSION VII41
